%% IEEE CEC 2027 submission draft (8 pages including references, IEEE conference format, double-anonymous review).
%% Scope: identical parallel machines (P||Cmax) and the permutation flowshop (RQ6). The unrelated-machine benchmark, the literature audit in full, further offline
%% selectors (kNN, MLP) and the full pre-registration record are left for the journal version. Red text marks what is still to be decided.
%% Build: bash cec/build.sh   (draft)   |   bash cec/build.sh anon   (anonymous version for submission)
\documentclass[conference,a4paper]{IEEEtran}
\usepackage{cite}
\usepackage{amsmath,amssymb}
\usepackage{graphicx}
\usepackage{booktabs}
\usepackage{multirow}
\usepackage{algorithm}
\usepackage{algorithmic}
\usepackage{xcolor}
\usepackage{url}
\newif\ifdraftmarks \draftmarkstrue
\newif\ifanon \anonfalse % CEC 2027 is double-anonymous: build with `bash cec/build.sh anon` (names, e-mail, repository link and red notes are removed)
\newcommand{\todonote}[1]{\ifdraftmarks\textcolor{red}{\textbf{[#1]}}\fi}
\newcommand{\Cmax}{C_{\max}}
\hyphenation{learn-heu-ris-tics}

\begin{document}

\title{Budget-Fair Attribution of the Gains of Learning-Guided Iterated Greedy: Where Does Online Learning Pay?}

\ifanon
\author{\IEEEauthorblockN{Anonymous Authors}}
\else
\author{\IEEEauthorblockN{Carlos Vinicio Dantas Ara\'ujo\IEEEauthorrefmark{1}\todonote{add advisor/coauthor}}
\IEEEauthorblockA{Graduate Program in Computer Science, State University of Cear\'a (UECE), Fortaleza, Brazil\\
\IEEEauthorrefmark{1}carlos.vinicio@aluno.uece.br}}
\fi

\maketitle

\begin{abstract}
Learning-guided metaheuristics are routinely reported to beat non-learning counterparts, yet comparisons seldom separate what the learned policy adds from what the search structure,
the CPU budget and the tuning already provide. We propose a budget-fair attribution protocol (a random control at a matched budget, equal CPU time with an efficient implementation of every arm,
an offline-validated competitor, oracle headroom, a power check and pre-recorded hypotheses) and apply it to an Iterated Greedy (IG) whose destruction is chosen by Q-Learning, bandits,
policy gradients, classic adaptive-operator-selection methods and neural learners, on identical parallel machines (synthetic classes and more than $13{,}000$ published instances) and on the permutation flowshop.
The destruction--reconstruction mechanism, not the policy, explains the advantage over Iterated Local Search, and the share of a local search's gain that is ``only budget'' ranges from $28\%$ to $93\%$ with the implementation.
On identical machines an online selector gains over random choice only where the oracle headroom is large; on the permutation flowshop it never does (headroom is necessary, not sufficient); and a static configuration (validated offline, or a uniform mixture) is never worse than the online learner.
The headroom can instead be exploited offline: a selector trained on $5{,}720$ disjoint instances captures about four fifths of it and beats both the validated action and the online learners, mostly through class-level knowledge.
A cheap pre-test on tuning instances ranks cells moderately well (AUC $0.73$--$0.75$ on a family never seen).
\end{abstract}

\begin{IEEEkeywords}
learnheuristics, adaptive operator selection, iterated greedy, parallel machine scheduling, permutation flowshop, experimental methodology
\end{IEEEkeywords}

\section{Introduction}
Metaheuristics that learn while they search are a prominent line of work in evolutionary computation and operations research. A typical \emph{learnheuristic}~\cite{CalvetEtal2017} embeds a
reinforcement-learning agent, most often tabular Q-Learning~\cite{WatkinsDayan1992,SuttonBarto2018}, as an adaptive selector of operators, strategies or parameters inside a search framework such as
Iterated Local Search (ILS)~\cite{LourencaoEtal2019} or Iterated Greedy (IG)~\cite{RuizStutzle2007}, and reports that it beats the search without the agent or published baselines~\cite{KarimiMamaghanEtal2023,MissaouiOSullivan2026}.
Such results are usually read as evidence that the \emph{learned policy} is what helps.

Three confounds make that reading fragile. First, the \emph{search structure} that hosts the agent (which operators, which acceptance rule, whether there is a local search) may produce most of the gain
under any selection rule. Second, the \emph{budget}: an agent changes the CPU time per iteration and a local search multiplies it, so equal iteration counts are not equal budgets. Third, the \emph{tuning}:
a parameter validated offline on separate instances is a competitor that most studies omit. Adaptive operator selection is a well-studied topic~\cite{FialhoEtal2010,PisingerRopke2007}, but we found it rarely
answered with these controls: in an audit of $23$ recent open-access papers that embed a learned selector in a metaheuristic, $10$ include a random-selection control, $6$ compare at a common time limit,
$3$ include a static configuration tuned offline and $22$ report a benefit of the learned component ($5$ of them only partly) (a small sample, one coder; details in the repository).

We ask six questions. \textbf{RQ1:} what explains the advantage of a learning-guided IG over ILS? \textbf{RQ2:} how much of the gain of a local search is budget, and does the answer depend on the implementation?
\textbf{RQ3:} do online learners beat random choice over the same actions, at equal iterations and at equal CPU time, does a static configuration win anyway, and could learning have shown up if it existed?
\textbf{RQ4:} where, if anywhere, does an online selector pay, and can that be foretold from tuning instances? \textbf{RQ5:} is the headroom exploitable at all, for instance by a policy learned offline?
\textbf{RQ6:} does any of this carry over to another problem? Our contributions are
(i) a \emph{budget-fair attribution protocol} (Section~\ref{sec:protocol}) with a power check, a positive control and a split of a learner's gain at equal CPU time into a decision effect and a cost effect, released with its code as a small library;
(ii) evidence on identical parallel machines that the advantage of the learning-guided IG comes from the search structure, and that the share of a local search's gain that is ``only budget'' depends on how efficiently the arms are implemented ($93\%$, $30\%$ and $28\%$ for loop-based, vectorized and compiled code);
(iii) a large-scale comparison with \emph{compiled} learners, a DQN and classic AOS methods in action spaces of $3$ to $30$ actions, at horizons of up to $10^5$ iterations and with finer times, whose answer is that a learner wins at equal CPU time mainly by choosing cheaper actions;
(iv) a \emph{headroom rule} for when an online selector pays, found on a first batch of published instances and confirmed on $3{,}950$ held-out ones, and a pre-test of it from tuning instances, validated out of sample;
(v) an \emph{offline} instance-conditioned selector that exploits the headroom the online learners leave, with controls that show where its gain comes from; and (vi) the same protocol on the permutation flowshop, with predictions registered beforehand, where no online learner gains and the headroom rule fails.

\section{Related Work}
The question is close to parameter control in evolutionary algorithms~\cite{EibenEtal1999,KarafotiasEtal2015,AletiMoser2016} and to selection hyper-heuristics~\cite{BurkeEtal2013,DrakeEtal2020}, which choose among low-level heuristics online;
offline configurators such as irace~\cite{LopezIbanezEtal2016} and ParamILS~\cite{HutterEtal2009} fix the parameters before the run, and tuning and control are two ways of configuring the same algorithm~\cite{EibenSmit2011}.
Machine learning inside metaheuristics has been surveyed~\cite{Talbi2021,KarimiMamaghanEtal2022}; where learned neighbourhood selection is reported to beat strong baselines, the policy is trained offline~\cite{SonneratEtal2021}, as in our selector (Section~\ref{sec:selector}).
Adaptive operator selection (AOS)~\cite{FialhoEtal2010} uses probability matching~\cite{Goldberg1990}, adaptive pursuit~\cite{Thierens2005}, ALNS weights~\cite{PisingerRopke2007}, bandits~\cite{AuerEtal2002,DaCostaEtal2008}
or reinforcement learning with a state~\cite{BengioBengio2021}. Q-Learning has been used to select operators inside an IG for the permutation flowshop~\cite{KarimiMamaghanEtal2023} and, closest to this work,
the destruction size of an IG for identical machines with setups~\cite{MissaouiOSullivan2026}; both report that the learning-guided method beats its baselines. That the component responsible for a gain
must be isolated by controlled experiments is a long-standing position~\cite{Hooker1995}; guidelines for comparing
optimization algorithms~\cite{LaTorreEtal2021,BartzBeielsteinEtal2020} and equivalence testing~\cite{Lakens2017} give the protocol below its footing.

\section{Method}
\subsection{Problem, instances and bounds}
In P$||\Cmax$, $n$ jobs with processing times $p_j$ are assigned to $m$ identical machines to minimise $\Cmax=\max_i\sum_{j\in S_i}p_j$; the problem is NP-hard~\cite{GareyJohnson1979}.
The classical bound is $\mathrm{LB}=\max(\max_jp_j,\lceil\sum_jp_j/m\rceil)$~\cite{Graham1969}, and $\mathrm{GAP}=100(\Cmax-\mathrm{LB})/\mathrm{LB}$. A tighter, certified bound comes from the linear relaxation of the arc-flow
model~\cite{MradSouayah2018}; it equals the best makespan found in $75.3\%$ of the synthetic instances (proving optimality). The synthetic instances have $p_j\sim U[1,100]$ in twelve classes
$(n,m)$ from $(20,4)$ to $(300,20)$ (and $p_j\sim U[1,10^4]$ in one regime below); each class has $100$ test instances and $30$ separate tuning instances. With integer times one unit of makespan is $0.12$--$0.78\%$ of it (median $0.38\%$), so a mean difference of
$0.05\%$ is an average over many ties and one-unit moves ($31$--$46\%$ of the instance pairs tie even after averaging five seeds). The published families B, C, E and F of the benchmark reviewed in~\cite{OstojicEtAl2025} give disjoint sets: $525$ tuning instances, $2{,}550$ \emph{discovery} and $3{,}950$ \emph{held-out}
test instances (families B, C, E), the $5{,}720$ B, C, E instances on which an offline selector is trained, and $780$ of family F (never used before); instances with processing times near $10^{18}$ do not fit the $64$-bit compiled core and were excluded ($378$ of the $7{,}403$ instances of the first three sets and $280$ of the $6{,}000$ training ones, all in family B).

\subsection{Search mechanisms and learners}
\label{sec:learners}
The \textbf{IG} starts from LPT (or from the best of LPT, MULTIFIT and LDM) and alternates destruction (removing $d$ jobs) and greedy reconstruction (non-increasing time onto the least-loaded machine), with greedy
acceptance $\Cmax(x')\leq\Cmax(x)$ and no separate perturbation phase. The \textbf{ILS} starts from LPT and, at each iteration, a selector chooses one of four neighbourhood operators (inter-machine swap,
reinsertion, intra-machine swap, block move) under best improvement; after $50$ iterations without improvement it restarts from the best solution after three random moves. The \emph{actions} are the destruction strategy (random, longest jobs, machine-concentrated; $3$ actions), the destruction size
($d/n\in\{0.05,0.10,0.20,0.30,0.40\}$; $5$ actions), their $15$ combinations, these $15$ with or without a descent after reconstruction ($30$ actions), or the four ILS operators.
The reward of an iteration is the improvement of the makespan, truncated at zero. At each iteration a selector $\pi$ chooses an action from the state of the current solution (the four features below), the IG destroys, reconstructs and optionally descends accordingly, $\pi$ receives the reward and the candidate is accepted if it does not worsen the makespan; the selector is the only component that changes between arms.

The \emph{learners} are: tabular Q-Learning over $81$ states ($4$ features $\times$ $3$ bins: gap to the bound, load imbalance, fraction of the budget used, stagnation; $\alpha=0.1$, $\gamma=0.9$,
$\varepsilon$-greedy from $1.0$ to $0.05$ by $0.995$ per step)~\cite{WatkinsDayan1992}; UCB1 with the reward normalised by its running maximum~\cite{AuerEtal2002}; a policy-gradient learner (softmax preferences per state,
a running-mean baseline, learning rate chosen on tuning instances)~\cite{Williams1992,SuttonBarto2018}; six methods without state, all tuned on the tuning instances: probability matching, adaptive pursuit, an ALNS-style roulette, Thompson sampling with discounting~\cite{Russo2018}, DMAB~\cite{DaCostaEtal2008} and an online racing learner (successive halving of the actions over the first half of the run); LinUCB on a continuous context~\cite{LiEtAl2010}; a one-hidden-layer neural Q-learner;
and a DQN with two layers of $64$ units, replay and a target network~\cite{MnihEtAl2015}. Random choice among the same actions is the control. The static competitor is the best fixed action per class,
chosen on the tuning instances.

\subsection{The budget-fair attribution protocol}
\label{sec:protocol}
(1)~\emph{Random control at a matched budget:} the effect of the policy is the difference between the learner and uniform random choice over the same actions, iterations and instances.
(2)~\emph{Equal CPU time, calibrated sequentially:} for each pair, the CPU time ratio $K$ is measured on separate instances, in sequence, with the process time and an otherwise quiet machine, and the cheaper
method receives $\mathrm{round}(1000K)$ iterations; as $K$ is a property of the implementation, every arm must be implemented with comparable efficiency, and $K$ is reported. A learner's gain at equal CPU time is the sum of a
\emph{decision effect}, its gain at equal iterations, and a \emph{cost effect}, the extra iterations the random control loses by matching the CPU time of a learner that chooses cheaper actions: $G_{\mathrm{CPU}}=G_{\mathrm{dec}}+G_{\mathrm{cost}}$ with
$G_{\mathrm{dec}}=-\bar d_{\mathrm{iter}}$ and $G_{\mathrm{cost}}=\bar d_{\mathrm{iter}}-\bar d_{\mathrm{CPU}}$.
(3)~\emph{A configuration validated offline} on separate instances is the competitor that a practitioner would use. (4)~\emph{Oracle headroom:} the gain of the best fixed action per instance over random
choice bounds what any selector that picks one action could gain. (5)~\emph{Power and a positive control:} the minimum detectable gain is estimated by planting gains in the paired differences (a fraction of the instances improved by one makespan unit),
and the learners and the pipeline are checked on a planted reward bonus and on a static policy chosen with class information. (6)~\emph{Statistics specified beforehand:} per instance (the mean of the seeds), $d=100(A-B)/B$ in the makespan is tested in pairs
(Wilcoxon, two-sided, ties dropped), with Holm correction over each family, a bootstrap $90\%$ interval, and an equivalence margin of $\pm0.05\%$, about an eighth of one makespan unit on the median synthetic instance; a Holm-significant difference whose interval lies inside
the margin is reported as parity. Throughout the text and tables a \emph{gain} is $-\bar d$ in percentage points of the makespan (written ``points'' or $\%$), positive when the first-named arm is better. Hypotheses, decision rules and predictions were recorded, with md5 hashes of the code, in a dated diary before each battery was run.

\subsection{Implementations}
The first implementation of the local search used Python loops. A second vectorized it (same moves, tie-breaking and results: $120$ of $120$ runs identical, $357$ unit tests). A third compiles the IG, the local
search, the destruction strategies and the agents (numba), with exact equivalence to the reference on the deterministic parts ($403$ and $35$ unit tests) and statistical equivalence on the random draws
(mean differences from $-0.003\%$ to $+0.012\%$ in five checks against the stored Python runs, all inside $\pm0.05\%$); it also counts evaluations.

\section{Results}
Table~\ref{tab:answers} summarizes the answers; the subsections give the evidence.

\begin{table}[t]
\caption{Answers to the six research questions (gains in percentage points of the makespan).}
\label{tab:answers}
\centering
{\scriptsize\setlength{\tabcolsep}{3pt}
\begin{tabular}{@{}lp{4.55cm}l@{}}
\toprule
 & Answer & Evidence \\
\midrule
RQ1 & The search structure: the IG beats the ILS by $0.46\%$ at equal CPU time; on the synthetic classes the policy adds at most $0.04$ points & \ref{sec:rq1} \\
RQ2 & It depends on the implementation: the share of a local search's gain that is budget is $93\%$, $30\%$ or $28\%$ & \ref{sec:rq2}, Table~\ref{tab:budget} \\
RQ3 & On the synthetic classes no learner beats random choice at equal iterations by more than $0.04$ points; equal-CPU wins are mostly a cost effect; a validated action wins & \ref{sec:rq3}, Table~\ref{tab:agents} \\
RQ4 & Where the headroom is large: Q-Learning gains $0.15$--$0.24$ points in family C and loses $0.05$--$0.07$ in family E, confirmed on held-out instances; a pre-test from tuning instances ranks the cells (AUC $0.73$--$0.75$ in an unseen family) & \ref{sec:published}, Tables~\ref{tab:published}, \ref{tab:pretest} \\
RQ5 & Yes, offline: a selector trained on $5{,}720$ instances beats the validated action by $0.02$--$0.04$ points and the online learners by $0.04$--$0.14$, mostly through class-level knowledge & \ref{sec:selector}, Table~\ref{tab:selector} \\
RQ6 & Partly: no online learner beats random choice, a validated action wins and an offline selector gains $0.03$--$0.05$ points, but headroom does not predict the online gain (rank correlation $-0.93$ and $-0.78$) & \ref{sec:flowshop}, Table~\ref{tab:flowshop} \\
\bottomrule
\end{tabular}}
\end{table}
\subsection{RQ1: the search structure, not the policy}
\label{sec:rq1}
Against LPT and seven ILS variants in the twelve classes, the Q-Learning-guided IG wins every one of the $72$ comparisons (of $96$) that remain significant after Holm correction, which is what a typical study would report as the success of the learned policy. Its advantage over the four ablation siblings
(random selection, UCB1, UCB1 per state, Q-Learning without state) is small: ten algorithm seeds on $240$ synthetic instances give a gain of $0.011\%$ ($90\%$ interval $0.003$ to $0.019$), real but negligible.
On the published instances the Q-Learning-guided IG beats random selection of the destruction strategy by $0.009$ to $0.026$ points
on $4{,}160$ held-out instances and by $0.015$ to $0.029$ on $2{,}690$ further ones (families B, C, E; Python implementation, all instances of those two sets; all six Holm-adjusted tests significant), and an informative state adds nothing
($-0.006$ points pooled, $p=0.77$). The search mechanism, by contrast, is large: with learning removed from both sides, the IG beats the ILS in $10$ of $12$ classes (the exceptions are the near-ceiling classes $(20,4)$ and $(40,5)$) by
$0.49\%$ on seed averages, against $0.53\%$ for the learned versions. This does not depend on how efficiently the ILS is implemented: with vectorized operators (identical results, $6.2$ times faster; an
iteration costing $0.4$ to $1.9$ times that of the IG), the IG still beats the ILS by $0.46\%$ at equal CPU time ($557$ wins against $87$) and by $0.50\%$ at equal iterations ($p_{\mathrm{Holm}}<10^{-80}$ in both), in all twelve classes at equal CPU time (from $0.004\%$ for $40\times5$ to $2.84\%$ for $60\times20$).

\subsection{RQ2: the budget attribution depends on the implementation}
\label{sec:rq2}
Adding a short best-improvement descent after each reconstruction improves the IG by $0.150\%$ at an equal number of iterations. The descent costs CPU time, so the verdict at equal CPU depends on the cost ratio $K$, and $K$ depends on the implementation (Table~\ref{tab:budget}).
With the first, loop-based descent ($K=9.6$), giving the plain IG the CPU time of the hybrid leaves an advantage of only $0.010\%$: $93\%$ of the gain would have been budget. Vectorizing the same descent cuts $K$ to $1.8$ and the advantage at equal CPU is $0.105\%$
($30\%$ budget); compiling both arms gives $K=2.0$ and $0.115\%$ ($28\%$), a reading registered before the compiled run (sixteen of sixteen predictions held). For the practical recipe (tuned $d$, start from the best constructive solution) the descent is neutral at equal CPU time in the vectorized and
compiled versions ($-0.008\%$ and $-0.006\%$, inside the margin). An evaluation count, proposed as an implementation-independent budget, agrees in sign but not in size ($0.080\%$, $50\%$ budget) and needs a stated cost model (weighting a reconstruction comparison by $0.25$ or $4$ moves the advantage between $0.026\%$ and $0.124\%$). Equal CPU time with an efficient implementation of every arm remains the budget of choice.

% Auto-generated by analysis/make_table_cec.py
\begin{table}[t]
\caption{Gain of adding a local search to the IG for P$||C_{\max}$ ($1{,}200$ synthetic instances per row), by budget convention and implementation. $K$: cost ratio of the hybrid to the plain IG (CPU time, or evaluation count). Gain in percentage points of the makespan (positive: the local search is better) with the 90\% interval; W/L/T: instances won, lost, tied by the hybrid. Share: the part of the equal-iterations gain that was budget.}
\label{tab:budget}
\centering
{\scriptsize\setlength{\tabcolsep}{1.7pt}
\begin{tabular}{@{}llrlrr@{}}
\toprule
Budget & Implementation & $K$ & $\bar d$ [90\% interval] & W/L/T & Share \\
\midrule
Equal iterations & Python & -- & +0.150 [+0.139, +0.162] & 436/16/748 & -- \\
Equal CPU & Python, loops & 9.6 & +0.010 [+0.002, +0.018] & 112/104/984 & 93\% \\
Equal CPU & Python, vectorized & 1.8 & +0.105 [+0.095, +0.115] & 330/18/852 & 30\% \\
Equal iterations & Compiled & -- & +0.159 [+0.147, +0.171] & 457/11/732 & -- \\
Equal CPU & Compiled & 2.0 & +0.115 [+0.104, +0.126] & 348/13/839 & 28\% \\
Equal evaluations & Compiled, count & 3.7 & +0.080 [+0.070, +0.089] & 267/27/906 & 50\% \\
\bottomrule
\end{tabular}}
\end{table}


\subsection{RQ3: learners, action spaces, the cost of the agent, and could learning have shown up?}
\label{sec:rq3}
Compilation makes a synthetic run about a hundred times faster, which allows five seeds per instance (the unit of analysis is the mean of the five, $1{,}200$ pairs per row) in four action spaces
(Table~\ref{tab:agents}). A compiled agent is cheap per unit of work: $1.0$--$1.09$ times random choice for UCB1 and the policy gradient and $1.07$--$1.33$ times for Q-Learning, which nevertheless needs
only $0.50$--$0.60$ of the CPU time of random choice in the three larger spaces because it learns to pick the cheaper actions. \emph{At an equal number of iterations no learner
gains more than $0.016\%$ over random choice in any of the $12$ comparisons}, within the equivalence margin.
\emph{At equal CPU time Q-Learning does win} in the two spaces of sizes, by $0.059\%$ and $0.049\%$, \emph{but only by choosing cheaper actions}: its decision effect is nil ($-0.003$ and $-0.004$ points) and its cost effect is $0.062$ and $0.053$. UCB1 and the policy gradient
gain nothing in either convention ($|\bar d|\leq0.008\%$ at equal CPU). Among the classic AOS methods the ALNS-style roulette is the exception: it gains $0.039$ points at equal iterations with five sizes
(a decision effect, Holm-significant and inside the margin; the cost effect adds $0.029$, so $0.068$ at equal CPU time, outside it) and nothing with fifteen actions ($-0.002$); probability matching and adaptive pursuit stay within $0.01$ points, Thompson sampling and DMAB within $0.013$, and online racing, the within-run analogue of a racing configurator, is \emph{worse} than random choice ($+0.012$ and $+0.067$ points).
The configuration validated per class beats every one of the nine methods by $0.053$--$0.147$ points in the spaces of $5$ and $15$ actions ($p_{\mathrm{Holm}}<10^{-57}$)
and is level with the three original learners in the other two, where the oracle headroom is smaller. Across the $48$ cells (class by space) the headroom does not predict the gain of Q-Learning (Spearman $\rho=0.09$, Fig.~\ref{fig:class}).
Where the actions differ in effect and cost (the four ILS operators, $12$ classes, $1{,}200$ instances) a DQN with two layers of $64$ units and LinUCB lose to the uniform mixture ($0.041\%$ and $0.066\%$), as does the operator validated per class ($0.062\%$);
over $15$ strategy--size combinations the DQN is level with random choice ($+0.023\%$) and the validated action beats it by $0.094\%$.

\emph{Could learning have shown up?}\label{sec:longnote} \emph{Power.} With $1{,}200$ pairs the analytic minimum detectable gain at $80\%$ power is $0.006$--$0.014$ points; the rank test under Holm correction is less sensitive, but planting gains in the paired differences gives, in the ten comparisons with a negligible observed effect, $58$--$99.8\%$ power at $0.02$ points and at least $97\%$ at $0.03$,
while the largest gain of Q-Learning, UCB1 or the policy gradient is $0.016$: gains of $0.03$ points or more are excluded for them, not smaller ones. \emph{Positive controls.} The per-class validated action, a static policy chosen with class information, is detected at $0.092$ and $0.080$ points ($p<10^{-76}$), and with a reward bonus planted on one action
all six learners choose it at least $1.5$ times more often than uniform even for half a makespan unit ($2.5$ to $13.7$ times for a bonus of $5$). \emph{Horizon.} With five sizes and $T=5{,}000$ and $20{,}000$ the three learners stay within $0.005$ points of random choice
while the oracle headroom falls from $0.125\%$ at $T=1{,}000$ to $0.060\%$ and $0.036\%$; at $T=100{,}000$ they are still within $0.006$ points of it, the headroom is $0.015\%$, and the validated action, which reaches the lower bound in $92\%$ of the instances, still leads the learners by $0.012$--$0.017$ points ($p<10^{-12}$): the online policy never catches up with a static choice, but almost nothing is left to choose.
\emph{Finer times.} With $p_j\sim U[1,10^4]$ the validated action never reaches the lower bound and the oracle headroom is $0.142\%$ and $0.175\%$; no learner gains more than $0.015$ points at equal iterations, Q-Learning wins at equal CPU time by $0.059$ and $0.076$ points ($K=0.63$ and $0.62$: cheaper actions again), and the validated action beats every learner by $0.105$--$0.136$ points.

% Auto-generated by analysis/make_table_cec.py
\begin{table*}[t]
\caption{Compiled agents and further online methods in action spaces of $3$ to $30$ actions (P$||C_{\max}$; $12$ classes $\times$ $100$ instances, five seeds, unit = instance mean, $1{,}200$ pairs per row). Gains in percentage points of the makespan with the 90\% interval (positive: the first-named arm is better). Oracle: best fixed action per instance over random choice (one seed per action, hence biased upwards). $K$: CPU time over that of random choice for the same number of iterations (geometric mean of the per-class ratios; equal CPU uses each class's own ratio). Equal iterations: learner over random choice ($T=1000$); equal CPU: learner over random choice with $T=\mathrm{round}(1000K)$; validated: best fixed action per class (tuning instances) over the learner, at equal iterations. The other five methods (probability matching, adaptive pursuit, Thompson sampling, DMAB, racing) are in the text.}
\label{tab:agents}
\centering
{\scriptsize\setlength{\tabcolsep}{3pt}
\begin{tabular}{@{}lllrlll@{}}
\toprule
Actions & Oracle & Learner & $K$ & Equal iterations & Equal CPU & Validated vs learner \\
\midrule
\multirow{3}{*}{3 strategies} & \multirow{3}{*}{0.046} & Q-Learning & 1.07 & +0.016 [+0.011, +0.021] & +0.011 [+0.006, +0.016] & -0.016 [-0.023, -0.010] \\
 &  & UCB1 & 1.01 & +0.002 [-0.003, +0.007] & -0.000 [-0.005, +0.005] & -0.002 [-0.009, +0.004] \\
 &  & Policy grad. & 1.01 & -0.003 [-0.008, +0.002] & -0.005 [-0.010, -0.000] & +0.003 [-0.003, +0.009] \\ \midrule
\multirow{4}{*}{5 sizes} & \multirow{4}{*}{0.125} & Q-Learning & 0.60 & -0.003 [-0.011, +0.005] & +0.059 [+0.050, +0.068] & +0.095 [+0.087, +0.102] \\
 &  & UCB1 & 1.01 & +0.001 [-0.005, +0.006] & -0.006 [-0.011, -0.001] & +0.091 [+0.085, +0.097] \\
 &  & Policy grad. & 1.03 & -0.000 [-0.005, +0.005] & -0.005 [-0.010, +0.000] & +0.092 [+0.086, +0.098] \\ \cmidrule(l){3-7}
 &  & ALNS roulette & 0.81 & +0.039 [+0.033, +0.045] & +0.068 [+0.062, +0.075] & +0.053 [+0.048, +0.058] \\ \midrule
\multirow{4}{*}{15 (strat.$\times$size)} & \multirow{4}{*}{0.122} & Q-Learning & 0.59 & -0.004 [-0.012, +0.004] & +0.049 [+0.040, +0.058] & +0.084 [+0.076, +0.092] \\
 &  & UCB1 & 1.03 & +0.004 [-0.001, +0.009] & +0.001 [-0.004, +0.005] & +0.076 [+0.070, +0.082] \\
 &  & Policy grad. & 1.06 & +0.000 [-0.005, +0.005] & -0.008 [-0.013, -0.003] & +0.080 [+0.074, +0.086] \\ \cmidrule(l){3-7}
 &  & ALNS roulette & 0.78 & -0.002 [-0.008, +0.003] & +0.020 [+0.014, +0.025] & +0.082 [+0.076, +0.089] \\ \midrule
\multirow{3}{*}{30 (+ local search)} & \multirow{3}{*}{0.049} & Q-Learning & 0.50 & -0.017 [-0.023, -0.011] & +0.016 [+0.010, +0.022] & +0.029 [+0.022, +0.035] \\
 &  & UCB1 & 1.05 & +0.004 [+0.000, +0.007] & +0.002 [-0.001, +0.005] & +0.008 [+0.004, +0.012] \\
 &  & Policy grad. & 1.08 & -0.001 [-0.005, +0.003] & -0.002 [-0.006, +0.001] & +0.013 [+0.008, +0.017] \\
\bottomrule
\end{tabular}}
\end{table*}


\begin{figure}[t]
\centering
\includegraphics[width=\columnwidth]{fig_cec4.pdf}
\caption{Oracle headroom against the gain of Q-Learning over random choice at equal iterations, per cell: synthetic class by action space ($48$ cells, grey), published family by size tercile by space ($18$ cells in the discovery batch, hollow, and $18$ in the held-out batch, filled), and flowshop class by action space ($18$ cells, crosses; $\rho$ per space). The oracle uses one seed per action and is therefore biased upward by run-to-run noise; the dashed line is gain equal to headroom.}
\label{fig:class}
\end{figure}
\subsection{RQ4: where does an online selector pay, and can that be foretold?}
\label{sec:published}
The synthetic classes have little headroom (pooled over classes the per-instance oracle gains at most $0.13\%$ over random choice, single cells up to $0.28\%$). The published benchmark~\cite{OstojicEtAl2025} is heterogeneous: families B, C and E with $4$ to $261$ jobs.
We ran the nine learners (compiled, five seeds, $T=1000$, five and fifteen actions) on $2{,}550$ of its instances, with the hyperparameters and the validated action fixed per family and size tercile on $525$ separate ones. \emph{Here learners do beat random choice} (Table~\ref{tab:published}, discovery):
Q-Learning gains $0.058$ points at equal iterations with fifteen actions and the ALNS roulette $0.057$ with five sizes, both outside the margin, and at equal CPU time Q-Learning gains $0.121$ and $0.137$ points. The gains follow the headroom (Fig.~\ref{fig:class}):
in family C, middle and upper size terciles (headroom $0.24$--$0.34\%$), Q-Learning gains $0.16$ and $0.24$ points, four fifths to nine tenths of the validated action's gain; in family E (headroom $0.07$--$0.11\%$) it \emph{loses} $0.06$--$0.07$ points,
and it loses in the smallest tercile of family B. The rank correlation between headroom and gain over the $18$ cells (family by tercile by space; the two spaces share instances) is $0.53$ ($0.37$ and $0.57$ within each space), against $0.09$ in the synthetic cells.
This was an exploratory reading of a batch whose registered predictions had missed in four of nine cases, so we fixed it as hypotheses before running $3{,}950$ held-out instances, retuning nothing. All five confirmatory tests held (Holm over five):
Q-Learning gains $0.245$ and $0.149$ points in family C, middle and upper terciles, with fifteen and five actions, and loses $0.052$ and $0.059$ in family E; the ALNS roulette gains $0.056$ with five sizes; the rank correlation was $0.65$ ($0.67$ and $0.60$ within each space, nine cells each, $p=0.05$ and $0.09$);
eleven of the twelve registered predictions held. The static competitor remains the safe choice: the validated action beats Q-Learning in all comparisons (by $0.064$--$0.082$ points in family C, middle and upper terciles) and every other learner,
and is level with the ALNS roulette only with five sizes ($0.021$ points, inside the margin).
Given the CPU time of each learner (the cost ratio of the validated action measured on its own), the validated action still beats all nine learners in all $54$ learner--space--set pairs (synthetic, discovery and held-out; Holm): by $0.071$--$0.097$ points against Q-Learning on the synthetic instances, $0.051$ and $0.058$ on the discovery and $0.058$ and $0.059$ on the held-out ones, and it is level with the ALNS roulette with five sizes on the published instances ($0.023$ and $0.024$).

% Auto-generated by analysis/make_table_cec.py
\begin{table}[t]
\caption{Published instances (compiled agents, five seeds, $T=1000$): gain over random choice in percentage points (positive: better); validated action: best fixed action per family and size tercile on separate instances. All: 2,550 (discovery) and 3,950 (held-out) instances; C$_{1-2}$: family C, middle and upper size terciles (700 and 1,020); E: family E (660 and 1,160). The held-out set was run after fixing the hypotheses (Section~\ref{sec:published}), retuning nothing.}
\label{tab:published}
\centering
{\scriptsize\setlength{\tabcolsep}{2.6pt}
\begin{tabular}{@{}lrrrrrr@{}}
\toprule
 & \multicolumn{2}{c}{All} & \multicolumn{2}{c}{C$_{1-2}$} & \multicolumn{2}{c}{E} \\ \cmidrule(lr){2-3}\cmidrule(lr){4-5}\cmidrule(l){6-7}
 & S5 & S15 & S5 & S15 & S5 & S15 \\
\midrule
\multicolumn{7}{@{}l}{\emph{Discovery}} \\
Oracle headroom & +0.120 & +0.169 & +0.243 & +0.336 & +0.067 & +0.107 \\
Q-Learning, equal iterations & +0.021 & +0.058 & +0.158 & +0.244 & -0.068 & -0.058 \\
ALNS roulette, equal iterations & +0.057 & +0.014 & +0.178 & +0.082 & -0.010 & -0.016 \\
Q-Learning, equal CPU time & +0.121 & +0.137 & +0.408 & +0.426 & -0.043 & -0.033 \\
Validated action & +0.075 & +0.112 & +0.194 & +0.273 & +0.027 & +0.038 \\ \midrule
\multicolumn{7}{@{}l}{\emph{Held-out}} \\
Oracle headroom & +0.132 & +0.174 & +0.305 & +0.378 & +0.068 & +0.105 \\
Q-Learning, equal iterations & +0.014 & +0.050 & +0.149 & +0.245 & -0.059 & -0.052 \\
ALNS roulette, equal iterations & +0.056 & +0.006 & +0.197 & +0.054 & -0.010 & -0.018 \\
Q-Learning, equal CPU time & +0.124 & +0.140 & +0.465 & +0.494 & -0.036 & -0.026 \\
Validated action & +0.078 & +0.101 & +0.231 & +0.308 & +0.022 & +0.036 \\
\bottomrule
\end{tabular}}
\end{table}


\emph{A pre-test.}\label{sec:pretest} The oracle headroom needs one run per action on every instance. A cheaper pre-test uses tuning instances only: the \emph{validated-gain-to-noise ratio} (VNR) is the gain of the best fixed action over random choice, cross-fitted on two halves of the tuning instances, divided by the run-to-run standard deviation of random choice
(relative, in \%, over three seeds), and the rule is ``an online selector pays in this cell if VNR $\geq\tau$''. We fixed $\tau=0.877$ on $42$ synthetic and discovery cells (AUC $0.78$, accuracy $0.83$) and applied it, unchanged, to the $26$ classes (subfolder, $m$, $n$) of family F, which no learner, tuning or threshold had seen:
$12$ tuning and $18$ test instances per class, $21$ cells per action space with a defined VNR (in the other five every run sits at the lower bound). The AUC is $0.75$ and $0.73$, the accuracy $0.71$ and $0.67$ and the rank correlation with Q-Learning's gain $0.65$ and $0.348$ in the spaces of five and fifteen actions (Table~\ref{tab:pretest}).
Nine of the $14$ registered predictions held; the miss that matters is the rank correlation in the larger space ($0.348$ against a registered $0.35$), so the pre-test is a useful ranking, not a decision rule, in a family whose effects are mostly below $0.1$ points. (There, Q-Learning gains $0.017$ and $0.016$ points, UCB1 $0.004$ and ALNS $0.026$ in the smaller space and $-0.003$ in the larger; the validated action is level or better in all six comparisons.)

% Auto-generated by analysis/make_table_cec2.py
\begin{table}[t]
\caption{Pre-test of the headroom rule: the validated-gain-to-noise ratio (VNR, estimated on tuning instances only) against whether Q-Learning beats random choice in a cell, with the threshold $\tau=0.877$ fixed on the first row. Cells with zero noise (every run at the lower bound; $5$ of $26$ per space in family F, $4$ of $18$ in the synthetic flowshop) are excluded. Accuracy: the rule `VNR $\geq\tau$' predicts `Q-Learning gains'; AUC is undefined (n/a) when no evaluable cell has a gain.}
\label{tab:pretest}
\centering
{\scriptsize\setlength{\tabcolsep}{3pt}
\begin{tabular}{@{}lrrrrr@{}}
\toprule
Cells & $n$ & Q gains & AUC & Acc. & Spearman \\
\midrule
P$||C_{\max}$ calibration cells (fixes $\tau$) & 42 & 0.45 & 0.78 & 0.83 & +0.589 \\
P$||C_{\max}$ held-out cells (explored earlier) & 18 & 0.44 & 0.93 & 0.89 & +0.843 \\
Family F, S5 (out of sample) & 21 & 0.62 & 0.75 & 0.71 & +0.655 \\
Family F, S15 (out of sample) & 21 & 0.62 & 0.73 & 0.67 & +0.348 \\
Flowshop, synthetic (out of sample) & 14 & 0.00 & n/a & 0.93 & -0.209 \\
Flowshop, Taillard (out of sample) & 14 & 0.07 & 0.23 & 0.86 & -0.407 \\
\bottomrule
\end{tabular}}
\end{table}


\subsection{RQ5: is the headroom exploitable at all?}
\label{sec:selector}
If online learners cannot capture the headroom, a policy learned \emph{offline} might. We trained an ExtraTrees regressor on $5{,}720$ B, C, E instances disjoint from every tuning and test set, with scale-free features of the instance ($n$, $m$, dispersion and quantiles of the times, bound ratios, the gap of LPT; no family label), to predict the relative performance of each fixed action; at test time it picks one action per instance.
The model and the hypotheses were fixed before any test run. On the discovery and held-out instances (Table~\ref{tab:selector}) it beats the validated action by $0.024$--$0.044$ points ($p_{\mathrm{Holm}}<10^{-38}$; inside the margin except in the larger space on held-out instances), also when the validated action is given the same number of evaluations ($0.027$--$0.039$; the selector runs $12$--$19\%$ more),
beats Q-Learning by $0.08$--$0.095$ and the ALNS roulette by $0.04$--$0.14$, and captures $0.79$--$0.84$ of the oracle headroom against $0.58$--$0.66$ for the validated action; all seven registered predictions held. Controls, registered after seeing this and run before looking at them, qualify the result: a static action learned on the same $5{,}720$ instances is only $0.001$--$0.019$ points better than the validated one, and a per-class lookup table learned on them
closes most of the remainder and is level with the selector on the held-out instances, where every class was seen in training ($0.004$ and $0.006$ points, not significant); on the discovery instances, where $59\%$ of them were, the selector wins by $0.014$ and $0.017$. Most of the gain is therefore knowledge about the class of instance, learned from more data, and the individual instance adds little.
Two further attempts at the online learner do not change the order: tuning Q-Learning per family ($27$ settings of learning rate, discount and exploration decay) adds $0.009$--$0.014$ points, and rewarding plateau moves, as the ALNS score does, adds $0.015$ with five sizes and \emph{costs} $0.007$--$0.009$ with fifteen; the validated action beats both in all eight comparisons and the selector by $0.06$--$0.10$.
The headroom is exploitable offline, from data, and not by adapting during a run.

% Auto-generated by analysis/make_table_cec2.py
\begin{table}[t]
\caption{Offline instance-conditioned selector (trained on $5{,}720$ disjoint instances, no family information) on the published instances ($2{,}550$ discovery, $3{,}950$ held-out; $T=1000$, five seeds): gain of the selector over each reference, in percentage points of the makespan (positive: the selector is better); at equal CPU time each arm has its own cost ratio $K$; $^\dagger$ not significant after Holm correction. Bottom: share of the oracle headroom captured (oracle of one seed per action, an upper bound; the cross-fitted oracle is a lower bound).}
\label{tab:selector}
\centering
{\scriptsize\setlength{\tabcolsep}{3pt}
\begin{tabular}{@{}lrrrr@{}}
\toprule
 & \multicolumn{2}{c}{S5} & \multicolumn{2}{c}{S15} \\ \cmidrule(lr){2-3}\cmidrule(l){4-5}
 & disc. & held & disc. & held \\
\midrule
vs.\ validated action ($525$ tuning inst.) & +0.024 & +0.026 & +0.027 & +0.044 \\
vs.\ validated action, equal CPU & +0.026 & +0.028 & +0.018 & +0.035 \\
vs.\ static action learned on the same $5{,}720$ & +0.022 & +0.025 & +0.018 & +0.025 \\
vs.\ per-class lookup table, same $5{,}720$ & +0.017 & +0.004$^\dagger$ & +0.014 & +0.006$^\dagger$ \\
vs.\ Q-Learning, equal iterations & +0.078 & +0.090 & +0.081 & +0.095 \\
vs.\ Q-Learning, equal CPU & +0.080 & +0.091 & +0.077 & +0.093 \\
vs.\ ALNS roulette, equal iterations & +0.042 & +0.048 & +0.125 & +0.140 \\
vs.\ ALNS roulette, equal CPU & +0.049 & +0.056 & +0.130 & +0.147 \\
vs.\ random choice & +0.099 & +0.104 & +0.139 & +0.146 \\
\midrule
Headroom captured: validated action & 0.62 & 0.59 & 0.66 & 0.58 \\
Headroom captured: lookup table & 0.69 & 0.75 & 0.74 & 0.80 \\
Headroom captured: selector & 0.83 & 0.79 & 0.82 & 0.84 \\
\bottomrule
\end{tabular}}
\end{table}


\subsection{RQ6: does it carry over? The permutation flowshop}
\label{sec:flowshop}
The same protocol ran on the permutation flowshop with the IG of Ruiz and St\"utzle (NEH start, Taillard's acceleration, insertion local search, constant-temperature acceptance) in a compiled core with the same agent code (reported: Q-Learning, UCB1, the policy gradient, the ALNS roulette and Thompson sampling; $39$ unit tests; destruction sizes $\{2,3,4,6,8\}$ and, in the larger space, three temperature factors; $T=1000$, three seeds).
It has nine classes ($n\in\{20,50,100\}$, $m\in\{5,10,20\}$) with $20$ tuning and $40$ test instances each, plus the $90$ instances ta001--ta090 of Taillard; tuning, validated actions, cost ratios and ten predictions were fixed beforehand. The classes with $m=5$ sit at the lower bound, and the headroom of the larger classes is $0.22$--$0.30\%$ (Table~\ref{tab:flowshop}).
\emph{No online learner beats random choice} at equal iterations: Q-Learning \emph{loses} $0.050$ and $0.061$ points on the synthetic instances and $0.059$ and $0.080$ on Taillard's (Holm-significant, not equivalent to zero) and still loses at equal CPU time ($0.036$ and $0.047$) although it costs $0.7$--$0.9$ of random choice; the other four learners stay within $0.013$ points on the synthetic instances.
The validated action beats the learners in all ten pairs at equal iterations (by $0.020$--$0.104$ points) and Q-Learning at equal CPU time (by $0.019$ and $0.064$). \emph{The headroom rule does not carry over:} the classes with the most headroom are those where Q-Learning loses most, and across the nine synthetic cells the rank correlation between headroom and Q-Learning's gain is $-0.93$ and $-0.78$ ($p<0.001$ and $0.013$).
Headroom is necessary, not sufficient: the run-to-run standard deviation of random choice is $0.2$--$0.3\%$ in the classes with $m=20$ and $n\geq50$, which swamps what an online selector can learn; the pre-test, which divides by that noise, stays below its threshold in $13$ of the $14$ evaluable cells but cannot discriminate, because no synthetic cell pays (one of the $14$ Taillard cells does, by $0.009$; Table~\ref{tab:pretest}). Four of the nine evaluable predictions held. An offline selector trained on $540$ further synthetic instances ($60$ per class, disjoint from the tuning and test instances; all six registered predictions held) gains $0.034$--$0.050$ points over random choice where Q-Learning loses, beats Q-Learning by $0.086$--$0.127$, and is level with a static action learned on the same instances (within $0.007$ points over all instances) and, not significantly, with the validated one: the headroom is exploitable offline and not online, at the level of the class.

% Auto-generated by analysis/make_table_cec2.py
\begin{table}[t]
\caption{Permutation flowshop (IG of Ruiz and St\"utzle, $T=1000$, three seeds): gain over random choice in percentage points of the makespan (positive: better) with five (F5) and fifteen (F15) actions. Synthetic: $360$ instances in nine classes $n\in\{20,50,100\}\times m\in\{5,10,20\}$; large: the four classes with $n\geq50$, $m\geq10$ ($160$); Taillard: ta001--ta090. Classes with $m=5$ sit at the lower bound.}
\label{tab:flowshop}
\centering
{\scriptsize\setlength{\tabcolsep}{1.8pt}
\begin{tabular}{@{}lrrrrrr@{}}
\toprule
 & \multicolumn{2}{c}{Synthetic} & \multicolumn{2}{c}{Synthetic, large} & \multicolumn{2}{c}{Taillard} \\ \cmidrule(lr){2-3}\cmidrule(lr){4-5}\cmidrule(l){6-7}
 & F5 & F15 & F5 & F15 & F5 & F15 \\
\midrule
Oracle headroom & +0.119 & +0.152 & +0.220 & +0.300 & +0.132 & +0.186 \\
Validated action & +0.028 & +0.043 & +0.036 & +0.092 & +0.012 & +0.031 \\
Q-Learning, equal iter. & -0.050 & -0.061 & -0.087 & -0.101 & -0.059 & -0.080 \\
ALNS roulette, equal iter. & -0.013 & -0.004 & -0.022 & -0.003 & -0.046 & -0.020 \\
Q-Learning, equal CPU time & -0.036 & -0.047 & -0.059 & -0.078 & -0.030 & -0.037 \\
Validated vs.\ Q-Learning, equal CPU & +0.019 & +0.064 & +0.011 & +0.114 & +0.013 & +0.080 \\
Static action ($540$ inst.) & +0.030 & +0.053 & +0.047 & +0.099 & +0.032 & +0.045 \\
Offline selector & +0.036 & +0.050 & +0.053 & +0.088 & +0.034 & +0.047 \\
\bottomrule
\end{tabular}}
\end{table}


\section{Discussion and Limitations}
The findings fit together. The gains that learning-guided search reports are mostly structure (destruction--reconstruction), budget (and how efficiently it is spent) and tuning. A learner's own contribution, the decision effect, is small where the oracle headroom is small, large where it is large (Q-Learning recovers most of a validated action's gain in family C of the published benchmark) and negative where headroom is small and the instances are heterogeneous (family E); where a learner wins at equal CPU time it does so largely through the cost of its actions.
What headroom does reward is knowledge acquired before the run: a selector trained offline on disjoint instances captures about four fifths of the oracle gain, mostly class-level knowledge that a lookup table also holds. The flowshop shows that headroom is necessary but not sufficient: where the run-to-run noise is high, the learner loses most where the headroom is largest.
\emph{A rule of use.} Estimate the oracle headroom, or its cheap proxy the VNR, on tuning instances: on identical machines an online selector is a liability below about $0.1$ points and can recover much of what a validated action gives above $0.2$, but on the flowshop headroom did not suffice; in all cases a validated action, or better a policy trained offline, is the baseline to beat.
A study that credits a learned selector should report the random control, the equal-CPU comparison with its cost ratio $K$, the competitor validated offline (or a uniform mixture), the oracle headroom, a power check with a positive control,
and the hypotheses fixed before the run with the predictions that failed; any adaptive component of a search can be audited the same way.

\emph{What the registered predictions showed.} The batteries were specified before they were run; $147$ of $183$ point predictions held in twenty batteries. The misses are informative: cost ratios of the agents, the headroom of an action space, tiny effects inside the margin, the classic AOS methods (the ALNS-style roulette gained, and won at equal CPU time, where we had predicted it would not),
the published instances (four of nine missed, which led to the headroom rule that the held-out replication confirmed) and the out-of-sample pre-test (a rank correlation $0.002$ below its threshold). The record also holds the corrections it forced on our text, the largest being the $93\%$ budget share, registered with the loop-based descent, that fell to $30\%$ once every arm was implemented with comparable efficiency.

\emph{Limitations.} (i) The equal-CPU comparison depends on $K$, measured on separate instances in sequence (not always on a quiet machine) and reported for every arm. (ii) Gains are in points of the makespan, so $0.05\%$ may or may not matter (one makespan unit is $0.12$--$0.78\%$ on the synthetic instances), and the margin is a choice.
(iii) Each battery has its own Holm family and registration, with no correction across batteries; the headroom rule began as an exploratory reading and was confirmed on instances that took no part in it, but its correlation rests on nine independent cells per space and is driven by the contrast between families C and E.
The oracle headroom is the best of one seed per action, so it is biased upward by run-to-run noise (the VNR is cross-fitted), and the learners start cold in every run, without a warm start across instances. (iv) The evidence is for P$||\Cmax$ and the permutation flowshop with greedy-acceptance IGs; $378$ published test instances with times near $10^{18}$ do not fit the $64$-bit compiled core and were excluded (all in family B); (v) Deep policy gradients and portfolios of hundreds of operators are outside the evidence, and the policy was studied up to $T=10^5$ with five sizes only. (vi) The pre-specification is a dated, hash-stamped diary, not a public registry; all verdicts, including the predictions that failed, are in the repository.
The study has not been externally reviewed. \todonote{decide: advisor as coauthor; Zenodo DOI for the record}

\section{Conclusion}
A learned operator selector is not what explains the advantage of an IG over ILS for identical parallel machines: that belongs to destruction--reconstruction. The share of a local search's gain that is budget is $93\%$, $30\%$ or $28\%$ depending on whether it is implemented with loops, vectorized or compiled, so the CPU cost ratio is itself a result to report.
The decision effect of an online selector follows the oracle headroom: nil to $0.04$ points on the synthetic classes, large where headroom is large (Q-Learning gains $0.15$--$0.24$ points in family C of the published benchmark, confirmed on $3{,}950$ held-out instances) and negative where it is small; a cheap pre-test from tuning instances ranks the cells moderately well.
Q-Learning wins at equal CPU time largely by choosing cheaper actions, and a static configuration was never worse than the online agent, while the tests can see gains of $0.03$ points and the learners capture a planted reward. The headroom the online learners leave is exploitable by a selector trained offline, though most of its edge is class-level knowledge. On the permutation flowshop no online learner beats random choice, the headroom rule fails (headroom is necessary, not sufficient) and a selector trained offline gains where the online learner loses.
Future work: deep policy-gradient learners, much larger operator portfolios, a compiled pipeline for unrelated machines, and the same protocol on population-based metaheuristics.

\vspace{2mm}
\noindent\textbf{Data and code availability.} Code, instances, results, the research diary with the pre-registered hypotheses and md5 hashes, the protocol library and the scripts that recompute every number are at
\ifanon an anonymized repository, whose link is given in the submission system.\else\url{https://github.com/CarlosAraujo7/pmsp}. \todonote{funding statement: CAPES/FUNCAP/CNPq to confirm}\fi

\bibliographystyle{IEEEtran}
\bibliography{refs_cec}

\end{document}
